\documentclass{ip-quiz}

\quizlevel{n1}
\quiznum{3}
\quiztitle{Funciones - Opción 1}

\begin{document}

\makeheader
\maketitle

\begin{sectionbox}{Enunciado}
Defina una función que reciba como parámetro una cantidad de segundos totales (un entero) y retorne un \texttt{str} con el tiempo equivalente en días, horas, minutos y segundos.

\textbf{Ejemplo:} Si la función recibe \texttt{93\_784}, debe retornar \texttt{"1 día(s), 2 h, 3 min y 4 s"}.
\end{sectionbox}

\begin{sectionbox}{Solución}
\begin{minted}{python}
def convertir_segundos(segundos_totales: int) -> str:
    segundos_en_un_dia = 60 * 60 * 24
    dias = segundos_totales // segundos_en_un_dia
    segundos_restantes = segundos_totales % segundos_en_un_dia

    segundos_en_una_hora = 60 * 60
    horas = segundos_restantes // segundos_en_una_hora
    segundos_restantes = segundos_restantes % segundos_en_una_hora

    segundos_en_un_minuto = 60
    minutos = segundos_restantes // segundos_en_un_minuto
    segundos = segundos_restantes % segundos_en_un_minuto

    return f"{dias} día(s), {horas} h, {minutos} min y {segundos} s"
\end{minted}
\end{sectionbox}

\end{document}
